\documentclass[12pt]{article}

\usepackage[margin=1in]{geometry}
\usepackage[utf8]{inputenc}
\usepackage[T1]{fontenc}
\usepackage{hyperref}
\usepackage{url}
\usepackage{booktabs}
\usepackage{amsfonts}
\usepackage{nicefrac}
\usepackage{microtype}
\usepackage{xcolor}
\usepackage{graphicx}
\usepackage{amsmath}
\usepackage{amssymb}

\title{Predicting Canadian Dollar Exchange Rates Against Major Global Currencies Using Machine Learning}
\author{Member 1 \and Member 2 \and Member 3}
\date{}

\begin{document}

\maketitle

\noindent Topic: Predicting Canadian Dollar Exchange Rates Against Major Global Currencies Using Machine Learning

\vspace{0.5em}
\noindent Group members: Member 1, Member 2, Member 3

\vspace{0.5em}
\noindent Abstract: Forecasting exchange rates is an important problem in finance and economics because currency movements influence international trade, investment decisions, and monetary policy. In this project, we study the behavior and prediction of the Canadian dollar (CAD) exchange rate relative to several major global currencies, including the US dollar (USD), euro (EUR), and Chinese Yuan (CNY), using daily exchange-rate data from the Bank of Canada. This problem is challenging because exchange-rate data are noisy, time-dependent, and influenced by many interacting factors. These properties make it a suitable machine learning task, especially from a statistical learning perspective.

Our motivation is to investigate whether supervised learning methods can identify predictive structure in historical exchange-rate data and provide better forecasts than conventional baseline methods. We formulate the task as a regression problem in which past daily exchange-rate observations are used to predict future values over a selected forecasting horizon. Our approach will begin with data preprocessing, followed by feature construction using lagged observations and rolling averages. We compare several machine learning models, including linear regression as a benchmark, support vector regression (SVR) to capture nonlinear relationships, and neural network models such as multilayer perceptrons (MLP). Model performance will be evaluated using standard error metrics such as mean absolute error (MAE) and root mean squared error (RMSE). Through this project, we aim to assess the effectiveness and limitations of machine learning methods for Canadian Dollar exchange-rate forecasting and better understand how statistical learning tools can be applied in financial time-series prediction.

\end{document}